\documentclass{tufte-handout}

\title{Cubic Reciprocity for \(x^3-2\)}

\author[Asher]{Asher Lucas-Cuddeback}

%\geometry{showframe} % display margins for debugging page layout

\usepackage{graphicx} % allow embedded images
\setkeys{Gin}{width=\linewidth,totalheight=\textheight,keepaspectratio}
\graphicspath{{graphics/}} % set of paths to search for images
\usepackage{amsmath}  % extended mathematics
\usepackage{amsthm}
\usepackage{booktabs} % book-quality tables
\usepackage{units}    % non-stacked fractions and better unit spacing
\usepackage{multicol} % multiple column layout facilities
\usepackage{lipsum}   % filler text
\usepackage{fancyvrb} % extended verbatim environments
\fvset{fontsize=\normalsize}% default font size for fancy-verbatim environments

% Standardize command font styles and environments
\newcommand{\doccmd}[1]{\texttt{\textbackslash#1}}% command name --
% adds backslash automatically
\newcommand{\docopt}[1]{\ensuremath{\langle}\textrm{\textit{#1}}\ensuremath{\rangle}}%
% optional command argument
\newcommand{\docarg}[1]{\textrm{\textit{#1}}}% (required) command argument
\newcommand{\docenv}[1]{\textsf{#1}}% environment name
\newcommand{\docpkg}[1]{\texttt{#1}}% package name
\newcommand{\doccls}[1]{\texttt{#1}}% document class name
\newcommand{\docclsopt}[1]{\texttt{#1}}% document class option name
\newenvironment{docspec}{
\begin{quote}\noindent}{
\end{quote}}% command specification environment
% Galois group: \Gal{L/K} or \Gal{L}
\newcommand{\Gal}[1]{\operatorname{Gal}\left(#1\right)}
\newtheorem{thm}{Theorem}
\begin{document}

\maketitle

\begin{thm}[{\citep[Theorem 4.15]{cox_primes_2022}}]
  The polynomial \(x^3-2\) splits modulo \(p\) if and only if
  \(p=a^2+27b^2\) for integers \(a\) and \(b\).
\end{thm}
\begin{proof}
  \textbf{(\(\Rightarrow\))} Suppose \(x^3-2\) splits modulo \(p\).
  If \(\lambda_1, \lambda_2\) are distinct roots in \(\mathbb{F}_p\)
  then \(\lambda_1 / \lambda_2 \in\mathbb{Z}[\omega]\)\sidenote{Since
    if \(\lambda_1\) is a root of \(x^3-2\) then
    \(\lambda_1^3=2\). Thus, \((\lambda_1 / \lambda_2)^3=2/2=1\) giving
  \(\lambda_1 / \lambda_2\in\mathbb{Z}[\omega]\)}.

  This implies that the unit group, \(\mathbb{F}_p^\times\), contains
  an element of order 3. Which by a theorem from Fermat gives us that

  \[3\,|\,p-1 \implies p\equiv 1\pmod{3}.\]

  With the original assumption that 2 is a cube modulo \(p\) we can state

  \[x^3-2 \text{ splits modulo } p\iff p\equiv 1\pmod{3}\text{ and 2
  is a cube modulo } p.\]

  Our prime \(p\) has this first condition \(p\equiv 1\pmod{3}\) if
  and only if \(p\) splits into primes in \(\mathbb{Z}[\omega]\) with
  \(p=\pi\overline{\pi}\), where \(\pi=a+b\omega\) and \(a\equiv 1\pmod{3}\).

  \newthought{Recall that} an Eisenstein integer is prime when, if
  \(a,b\neq 0\), \(||\pi||= a^2-ab+b^2\) is an ordinary prime in \(\mathbb{Z}\).

  We can clean up this norm by exploiting the fact that we are in
  \(\mathbb{Z}[\omega]\) to try and eliminate the \(-ab\) term.
  Here, we can, instead of using \(\pi=a+b\omega\), use
  \(\omega\pi=a\omega+b\omega^2=a\omega+b(-1-\omega)\) which leaves
  the norm unchanged.
  \sidenote{The norm is multiplicative so \(||\omega\pi||=
  ||\omega||\cdot||\pi||=(\omega\overline{\omega})||\pi||=(\omega\cdot\omega^2)||\pi||=||\pi||\)}

  This transposition continues with:
  \[a\omega+b(-1-\omega)=-b + (a-b)\omega\]
  Whereupon we rename \(A=-b\) and \(B=(a-b)\) to re-obtain
  \(A+B\omega\). This process is the same for \(\omega^2\pi\) where we get
  \(A=b-a\) and \(B=-a\). Since \(p=\pi\overline{\pi}\) is an odd
  prime, \(a\) and \(b\) cannot both be even.

  Making this substitution against the norm equation and using
  the reformulation that yields and even value for \( b \) allows

  \[
    ||\pi || = a^2 -ab + b^2 = a^2 - 2aB + 4B^2
  \]
  Which collapses to, after completing the square
  \[
    = (a-B)^2 + 3B^2 = A^2 + 3B^2
  \]

  \[||\pi||=a^2+3b^2\]
  Using the theorem\sidenote{I don't know where this comes from.}
  which states
  \[2\text{ is a cube modulo }p\iff b\equiv 0\pmod{3}\]
  it becomes obvious that \(b=3b_0\). So, by simply plugging this
  value in, we find:
  \[p=a^2+3(3b_0)^2=a^2+27b^2\]
  as desired.
  \newline
\end{proof}

The above proof is an \emph{attempted} proof by amalgamation of
observing the proof of this theorem from the time before class field
theory. The proof that follows is attempt to appeal even more to Galois groups.

\begin{proof}
  The polynomial \( p(x) = x^3 - 2 \), via the fundamental theorem of
  algebra, has three roots. Because \( p \) is a binomial equation,
  we know that all the roots must exist as cyclic permutations of
  one root in \( \mathbb{C} \). The roots of \( p \) are:
  \[
    \{ \sqrt[3]{2}, \omega\sqrt[3]{2}, \omega^2\sqrt[3]{2} \}
  \]
  where \( \omega \) is the third root of unity\sidenote{\(
  \omega^3=1 \)}. The splitting field of \( p \) is plainly, then, \(
  \mathbb{Q}(\omega, \sqrt[3]{2}) \). Examining the tower of field
  extensions gives
  \[
    K=\mathbb{Q}(\omega)
  \]
  as a degree two and
  \[
    K(\sqrt[3]{2})
  \]
  as a degree three extension. Thus \( \Gal{\mathbb{Q}(\omega,
  \sqrt[3]{2}) / \mathbb{Q}} \) is a group of order 6 by the splitting
  correspondence. The automorphisms which generate \(
  \Gal{\mathbb{Q}(\omega, \sqrt[3]{2})} \) are
  \[
    \begin{matrix}
      \text{fix }\omega & \lambda \mapsto \lambda\omega \\
      \text{fix }\lambda & \omega\mapsto\omega^2
    \end{matrix}
  \]
  All of these facts add up to tell us that \(
  \Gal{\mathbb{Q}(\omega, \sqrt[3]{2}) / \mathbb{Q}} \cong S_3 \)
  which is a solvable group.
\end{proof}
\bibliographystyle{plainnat}
\bibliography{cubic-reci}
\end{document}
